% BHU PYQ -- Auto-generated & error-corrected
\documentclass[11pt,a4paper]{article}

\usepackage[a4paper,top=2.5cm,bottom=2.5cm,left=2.8cm,right=2.8cm,headheight=14pt]{geometry}
\usepackage{amsmath,amssymb}
\usepackage[T1]{fontenc}
\usepackage[utf8]{inputenc}
\usepackage{lmodern}
\usepackage{microtype}
\usepackage{fancyhdr}
\usepackage{enumitem}
\usepackage{hyperref}
\usepackage{lastpage}
\usepackage{parskip}

\hypersetup{colorlinks=true,urlcolor=black,linkcolor=black,
  pdftitle={STA-101: Statistical Methods and Probability -- BHU PYQ},pdfauthor={Banaras Hindu University}}

\pagestyle{fancy}\fancyhf{}
\renewcommand{\headrulewidth}{0.4pt}\renewcommand{\footrulewidth}{0.4pt}
\fancyhead[L]{\small   \textit{Banaras Hindu University}}
\fancyhead[R]{\small   \textit{STA-101: Statistical Methods and Probability}}
\fancyfoot[C]{\small Page  \thepage\ of \pageref{LastPage}}


\newlist{parts}{enumerate}{2}
\setlist[parts,1]{label=(\alph*),leftmargin=*,itemsep=5pt,topsep=3pt}
\setlist[parts,2]{label=(\roman*),leftmargin=*,itemsep=3pt,topsep=2pt}

\newcommand{\pts}[1]{\hfill   \textit{[#1]}}


\begin{document}


\begin{center}
  {\large\bfseries BANARAS HINDU UNIVERSITY}\\[2pt]
  {\normalsize B.A. (Hons.) Semester I Examination, 2022-23}\\[6pt]
  \rule{0.8\linewidth}{0.5pt}\\[6pt]
  {\large\bfseries Paper No. STA-101: Statistical Methods and Probability}\\[4pt]
  {\normalsize Time: 3 Hours\hfill Maximum Marks: 70}
\end{center}
\rule{\linewidth}{0.4pt}
\smallskip



\noindent   \textit{Instructions: Answer any   \textbf{five} questions. All questions carry equal marks. Symbols have their usual meanings.}
\bigskip

\medskip


\noindent   \textbf{Question 1.}

\begin{parts}
  \item State the advantages of graphical representation of data. Explain the histogram with an example. \pts{7}
  \item Differentiate between primary data and secondary data. Describe the different sources of secondary data in India. \pts{7}
\end{parts}

\medskip\rule{\linewidth}{0.2pt}

\medskip


\noindent   \textbf{Question 2.}

\begin{parts}
  \item Describe the different measures of central tendency. Briefly explain the meaning of an outlier. Using an example, show how an outlier can affect the value of the arithmetic mean. \pts{7}
  \item Define median for ungrouped data when the number of observations is even. Derive the expression to obtain the median for a continuous frequency distribution. \pts{7}
\end{parts}

\medskip\rule{\linewidth}{0.2pt}

\medskip


\noindent   \textbf{Question 3.}

\begin{parts}
  \item State the relative and absolute measures of dispersion and describe the main difference between mean deviation and standard deviation. \pts{7}
  \item Define the relation between moments about the mean in terms of moments about an arbitrary point, and vice versa. \pts{7}
\end{parts}

\medskip\rule{\linewidth}{0.2pt}

\medskip


\noindent   \textbf{Question 4.}

\begin{parts}
  \item Describe the absolute measures of skewness. Explain the coefficient of skewness based upon moments ($eta_1$ and $\gamma_1$). \pts{7}
  \item Discuss the limitations of the classical definition of probability. For any two events $E_1$ and $E_2$, show that:
  \[
  P(E_1 \cup E_2) = P(E_1) + P(E_2) - P(E_1 \cap E_2).
  \]
  Also, for events $E_1, E_2, \ldots, E_n$, show that:
  
\begin{parts}
    \item $P(E_1 \cap E_2 \cap \cdots \cap E_n) \geq P(E_1) + P(E_2) + \cdots + P(E_n) - (n-1)$
    \item $P(E_1 \cup E_2 \cup \cdots \cup E_n) \leq P(E_1) + P(E_2) + \cdots + P(E_n)$
  \end{parts}\pts{7}
\end{parts}

\medskip\rule{\linewidth}{0.2pt}

\medskip


\noindent   \textbf{Question 5.}

\begin{parts}
  \item For any three events $E_1$, $E_2$ and $E_3$, show that:
  \[
  P(E_1 \cup E_2 \mid E_3) = P(E_1 \mid E_3) + P(E_2 \mid E_3) - P(E_1 \cap E_2 \mid E_3).
  \]
  \pts{7}
  \item In two companies A and B, a student of Statistics has applied for a job. The probability of his selection in company A is $0.7$ and of his rejection in company B is $0.5$. The probability of at least one of his applications being rejected is $0.6$. What is the probability that he will be selected in exactly one of the companies? \pts{7}
\end{parts}

\medskip\rule{\linewidth}{0.2pt}

\medskip


\noindent   \textbf{Question 6.}

\begin{parts}
  \item What do you understand by the distribution function of a random variable? Describe the important properties of a distribution function. \pts{7}
  \item What do you mean by a probability density function (p.d.f.)? Let $X$ be a continuous random variable with p.d.f.\ given by:
  \[
  f(x) = 
\begin{cases}
    px, & 0 < x < 1 \\
    p, & 1 \leq x < 2 \\
    -px + 3p, & 2 \leq x < 3 \\
    0, &  \text{otherwise}
  \end{cases}
  \]
  
\begin{parts}
    \item Determine the constant $p$.
    \item Determine the cumulative distribution function (CDF).
    \item If three independent observations are made, what is the probability that exactly one of these three numbers is larger than $1.5$?
  \end{parts}\pts{7}
\end{parts}

\medskip\rule{\linewidth}{0.2pt}

\medskip


\noindent   \textbf{Question 7.}

\begin{parts}
  \item Describe the conditional probability distribution function. The joint probability distribution of two random variables $X$ and $Y$ is given below:

  
\begin{center}
  
\begin{tabular}{c|cccc}
    \hline
    $X ackslash Y$ & $1$ & $2$ & $3$ & $4$ \\
    \hline
    $1$ & $1/9$ & $1/2$ & $1/18$ & $1/36$ \\
    $2$ & $1/2$ & $1/2$ & $1/8$ & $1/36$ \\
    $3$ & $5/36$ & $1/36$ & $1/36$ & $1/36$ \\
    $4$ & $1/36$ & $1/18$ & $1/36$ & $5/36$ \\
    \hline
  \end{tabular}
  \end{center}

  Find the conditional distribution of $X$ given $Y = 1$. \pts{7}
  \item Show that the mathematical expectation of the sum of $n$ random variables is equal to the sum of their individual expectations, provided all expectations exist. \pts{7}
\end{parts}

\vfill
\rule{\linewidth}{0.4pt}

\begin{center}\small
  \href{https://bhu-exam.vercel.app/}{Click here to visit our website}\\[4pt]
  \href{https://me-aryan.vercel.app/}{Click here to visit our contributor}\\[4pt]
  \href{https://quashan-me.vercel.app/}{Click here to visit our contributor}
\end{center}

\end{document}